\section{Discussion}\label{sec:discussion}

\paragraph{What the engine changes.} The baseline design keeps every gate on a partition qubit on the host and sweeps
the whole slice once per native operation. Its run time is therefore set by per gate kernel sweeps for wide circuits
and by synchronisation for small ones, as the model of Section~\ref{subsec:res-model} shows without any fitted
parameter. The engine removes both costs: gates that need no amplitudes of another slice run on the cards under
the slice enable (F1 to F3), consecutive records share one HBM pass, and the state uses all 32 pseudo channels of each
card. The QFT at $n=20$ on four cards takes \NumHostSubmittedQftTwenty~s in the baseline design and
\NumHostRevisedQftTwenty~s in the engine on the same cards, \NumSubmittedOverRevisedQftTwenty~times less.

\paragraph{Where the system stands.} We make no claim of speed advantage. Qiskit Aer with all threads of the host is
faster than the four card system for almost every circuit of Table~\ref{tab:cpu}; the FPGA system is faster than a
single CPU thread in all \NumCpuWideCases~cases with $n\ge 24$, and its time approaches that of the multithreaded CPU as
the width grows. Our measured HBM rate, \NumHbmPerCard~GB/s per card, is well below the peak of the memory, and the engine is compute bound
when a pass carries several operations; a design that issues more amplitude pairs per cycle, or more engines than the
eight that routed in complex128, would raise the rate of the cards. The comparison with GPUs on the same card type is in
the companion paper~\cite{tiwari2026fpga30}.

\paragraph{What distribution buys.} Each doubling of the cards adds one qubit of capacity in complex128, from 29 qubits on
one card to 31 on four (Section~\ref{subsec:res-maxwidth}; verified: \NumMaxWidthOneCard, \NumMaxWidthTwoCards~and
\NumMaxWidthFourCards~qubits), at the price of exchanges through host memory. Scaling and energy show the price:
more cards make most runs slower and every run more energy hungry (Sections~\ref{subsec:res-scaling}
and~\ref{subsec:res-energy}). The slice paths and the compiler rewrite keep the exchanges to a minimum where the circuit
allows it: without them the QFT on four cards is \NumAblQftFourAllOff~times slower and GHZ preparation
\NumAblGhzFourAllOff~times. Circuits whose slice qubits carry general gates, such as random circuits, gain nothing from
them and pay the full exchange cost.

\paragraph{Relation to other multi device simulators.} Distributed CPU and GPU simulators use the same global to local
swap and reach far larger widths on many nodes~\cite{haner2017petabyte,deraedt2019massively,imamura2022mpiqulacs}. The
Trefoil system distributes much larger states over FPGAs with SATA storage, trading bandwidth for
capacity~\cite{niwase2023trefoil,wei2025qutrefoil}. The present design keeps the state in HBM, which limits the width to
the memory of the cards but keeps every pass at HBM speed; its weak point is the exchange path between the cards, which
peer to peer transfers or direct card to card links would shorten.

\paragraph{Reproducibility.} Every number in this paper is generated by script from the raw measurement files, each of
which records the code commit, the bitstream hash and the cards used. The circuits, the transpilation settings, the
timing definitions and the measurement protocol are given in Section~\ref{sec:methodology}, and the raw data are
available as stated in the declarations.
